
\chapter{Подписки и~сохранённые реквизиты}\label{ch:subscriptions}
Подписка начинается с~первого согласия на~хранение реквизитов. \highlight{Для
продукта январский платёж и~февральское автосписание продолжают
одну услугу, но~для эмитента это две разные авторизационные
ситуации}. Когда связь между ними обрывается, эмитент теряет контекст,
карточная сеть считает операцию подозрительной, а~клиент получает отказ
по~штатному платежу.

\section{Сохранённые реквизиты и SCF}\label{sec:subscriptions-scf}

До~появления единых правил ТСП решали повторные списания каждый
по-своему: кто-то хранил номер карты (Primary Account Number, PAN),
кто-то полагался на~токенизацию, кто-то каждый месяц отправлял обычный
запрос и~надеялся, что эмитент угадает контекст. Главное эмитенту
не~сообщали: продолжение ли это согласованной цепочки или новая
операция без явного участия клиента.

Отсутствие контекста делало \enquote{тихую} транзакцию, инициированную
ТСП (merchant-initiated transaction, MIT), неотличимой от~обычного
запроса без присутствия карты (card-not-present, CNP). ТСП мог
начать регулярные списания без явного согласия клиента или продолжить
их после его отмены, и~эмитент не~имел технической возможности это
различить. SCF закрыл этот разрыв: контекст первого согласия стал
обязательным полем авторизационного сообщения.

\highlight{Visa и~Mastercard ввели Stored Credential Framework (SCF)~--- набор
правил, по~которым карточная сеть различает первое согласие клиента
и~последующие списания, делая контекст первого согласия видимым
в~авторизационном сообщении}~\cite{visa-scf,mc-tpr}. Реквизиты
остаются у~ТСП, но~каждая операция должна ясно показывать
связь с~исходным согласием клиента.

В~авторизационном сообщении эта связь выражается несколькими
признаками: первичным или последующим использованием сохранённых
реквизитов, типом сценария и~ссылкой на~исходную авторизацию,
связывающей всю цепочку~\cite{mc-tpr,cybs-initial-transactions}.
Visa использует Stored Credential Transaction Indicator со~значением
\texttt{initial} при первой транзакции, инициированной клиентом
(customer-initiated transaction, CIT), и~\texttt{subsequent} при
последующей MIT. Рядом передаётся тип инициатора
(cardholder-initiated или merchant-initiated) и~причина
MIT~--- регулярные (recurring), рассрочка (installment) или
внеплановые
(unscheduled)~\cite{visa-scf,cybs-initial-transactions}. Mastercard
кодирует ту же информацию через subelements поля данных~48 (Data
Element~48, DE~48; POS Entry Mode и~MIT Framework Indicator; точные
subfields сверяются по~действующей редакции правил обработки
транзакций, Transaction Processing Rules, TPR) и~требует передавать
Trace~ID исходной CIT в~каждой последующей MIT~\cite{mc-tpr}.

Клиент оформляет месячную подписку. Первый платёж идёт
как CIT с~indicator = \texttt{initial} и~получает от~карточной сети
идентификатор сетевой транзакции (Network Transaction~ID, NTID
у~Visa; Trace~ID у~Mastercard~--- функциональный аналог NTID
для соответствующей сети)~\cite{mrc-nti-trace-id}. Через
30~дней ТСП отправляет MIT с~indicator = \texttt{subsequent},
типом \texttt{recurring} и~ссылкой на~NTID первой CIT. Когда
ТСП не~передаёт ссылку на~исходную авторизацию или указывает
неверный тип, эмитент видит запрос без контекста и~вправе применить
стандартную политику CNP, вплоть до~отказа
(см.~рис.~\ref{fig:decline-decision-tree}). При аудите карточная сеть
квалифицирует такую операцию как нарушение правил
SCF~\cite{mc-tpr,visa-core-rules}.

\begin{figure}[htbp]
  \centering
  \includefiguresvg[width=\linewidth]{assets/figures/ch29-subscriptions/scf-lifecycle}
  \caption{SCF на~временной шкале: CIT с~выдачей NTID и~серия MIT
  со~ссылкой на~NTID; отсутствие NTID ведёт к~отказу.}\label{fig:scf-lifecycle}
\end{figure}
\FloatBarrier

\section{CIT и MIT}\label{sec:subscriptions-cit-mit}

Подписочный процесс ломается на~стыке двух событий. Клиент
оформил услугу и~согласился на~сохранение реквизитов, а~через
неделю или месяц ТСП пытается списать очередной платёж уже
без~его участия~--- и~контекст обрывается.

В~CIT клиент сам нажимает кнопку оплаты, вводит данные карты,
проходит аутентификацию, когда она требуется правилами рынка. Так
выглядит первый платёж за~подписку, оформление Card-on-File (CoF,
модель хранения реквизитов карты у~ТСП для последующих
списаний) или покупка, после которой ТСП получает право
сохранить реквизиты. Система фиксирует исходную точку цепочки
и~получает связующий идентификатор исходной авторизации.
В~документации сети или шлюза этот идентификатор называется
transaction identifier, network transaction ID,
\texttt{previousTransactionID} или Trace~ID~\cite{mc-tpr,cybs-initial-transactions,cybs-mc-subscription-mit}.

Согласие клиента возникает в~CIT, но~сохранение реквизитов
не~означает право на~любые будущие списания. Идентификатор нужно
хранить и~передавать в~каждой последующей операции, иначе связь
между первым согласием и~следующими списаниями исчезнет
из~авторизационного контекста~\cite{mc-tpr,cybs-initial-transactions}.

MIT~--- списание без активного участия клиента: ежемесячная
подписка, доначисление по~уже оказанной услуге или платёж
по~ранее согласованной схеме. Запрос инициирует ТСП,
но~обязан показать эмитенту, что операция продолжает существующую
цепочку. Поэтому MIT всегда сопровождается признаками SCF, типом
операции, заданным карточной сетью, и~идентификатором исходной
авторизации~\cite{mc-tpr,cybs-mc-subscription-mit}.

MIT без~корректных признаков SCF эмитент трактует как новую
подозрительную транзакцию и~отклоняет, даже когда ТСП считает
её штатной по~своей бизнес-логике. В~европейском трафике
при~удалённом оформлении мандата на~серию MIT (мандат~--- согласие
клиента на~регулярные списания по~сохранённым реквизитам) строгая
аутентификация клиента (strong customer authentication, SCA)
применяется на~этапе создания мандата, а~последующие MIT должны
быть корректно помечены и~связаны с~исходной
авторизацией~\cite{eba-mit-qa,mc-tpr}.

\section{Идентификатор исходной авторизации и~цепочка согласия}\label{sec:subscriptions-ntid}

\highlight{Токен отвечает на~вопрос, каким реквизитом платить, а~идентификатор
исходной авторизации~--- на~вопрос, с~каким ранее данным согласием
связано текущее автосписание}. Токен обновляется при перевыпуске
карты, а~идентификатор цепочки хранится отдельно от~самих
реквизитов. Mastercard в~Европейской экономической зоне (ЕЭЗ),
Великобритании и~Гибралтаре требует в~последующих запросах
авторизации по~регулярным платежам передавать уникальный Trace~ID
из~ответа на~исходную авторизацию, а~эмитент должен уметь по~нему
восстановить исходную транзакцию~\cite{mc-tpr}. У~Visa аналогичное
требование: после начальной операции нужно сохранить идентификатор
транзакции (\textit{transaction identifier}) для всех последующих
операций~\cite{cybs-initial-transactions}.

Visa возвращает NTID в~ответе на~авторизацию (\textit{Transaction
Identifier} в~DE~62.2, до~15~цифр, генерирует VisaNet); Mastercard
использует термин Trace~ID и~требует его передачу в~DE~48 SE~63
последующих авторизаций~--- 15~символов фиксированной структуры
FNC\,$\cdot$\,BRN\,$\cdot$\,MMDD\,$\cdot$\,spaces, с~2026~переименовывается
в~\textit{Transaction Link ID} (TLID)~\cite{mc-tpr,cybs-initial-transactions,mc-an11418-tlid}.
Trace~ID не~равен ARN (DE~31, 23~цифры). Индикаторы CoF/MIT передаются
у~Visa через SCF (C/R/I) и~POS~Environment DE~60.8; у~Mastercard~---
через DE~48 SE~22 (SF1, SF5); POS~Entry~Mode~10~--- единый код
\textit{Credentials on File} для обеих схем~\cite{visa-vbs-10,mc-spme-2025}.
У~Cybersource и~других шлюзов поле называется
\texttt{previousTransactionID} или
\texttt{networkTransactionID}~--- разные имена для одного
и~того же идентификатора.

Хранить идентификатор нужно на~уровне подписки. Каждый последующий
MIT ссылается на~одну и~ту же исходную CIT, поэтому идентификатор
живёт столько же, сколько сама подписка. Когда система хранит его
только в~записи последней транзакции, при любом сбое записи цепочка
обрывается. Поэтому в~модели подписки заводят отдельное поле,
которое заполняется при первой CIT и~не~перезаписывается
последующими ответами~\cite{cybs-mc-subscription-mit}.

При миграции между провайдерами платёжных услуг (payment service
provider, PSP) ТСП переносит токены и~историю подписок,
но~забывает перенести идентификатор исходной авторизации первой
CIT. На~следующем цикле MIT приходит уже без этой привязки,
и~ТСП теряет связующий идентификатор между первым согласием
и~последующим списанием~\cite{mc-tpr,cybs-mc-subscription-mit}.

\begin{figure}[htbp]
  \centering
  \includefiguresvg[width=\linewidth]{assets/figures/ch29-subscriptions/ntid-chain}
  \caption{Цепочка NTID/Trace~ID: форматы полей Visa и~Mastercard,
  индикаторы CoF/MIT, POS~Entry~Mode~10.}%
  \label{fig:ntid-chain}
\end{figure}

Если исходная CIT оформлена некорректно, поверхностные правки
цепочку согласия не~восстановят: идентификатор приходится
получать заново через новую CIT с~подтверждением клиента.

Связка CIT~$\to$~MIT работает не~только в~авторизации.
В~3DS рекуррентные операции идут как 3RI-запросы со~ссылкой
на~DS Transaction~ID исходной аутентификации
(см.~гл.~\ref{ch:3ds}, раздел про~3RI). В~спорах Visa Compelling
Evidence~3.0 (см.~раздел~\ref{sec:disputes-ce30}) та же ссылка
подтверждает, что текущее списание принадлежит уже знакомому
клиенту.

\section{Мягкий отказ и~повторная попытка}\label{sec:subscriptions-soft-decline}

Подробное разбиение мягких и~жёстких отказов и~стратегии повторов
вынесены в~главе~\ref{ch:decline-management}. Когда MIT получает
отказ, который не~закрывает сценарий окончательно, у~ТСП есть
несколько допустимых действий: новая аутентификация, повтор
в~разрешённый интервал, обработка кода рекомендаций ТСП
(Merchant Advice Code, MAC) или уведомление
клиента~\cite{visa-core-rules,mc-tpr}.

Когда отказ требует аутентификации, клиента возвращают в~новый
подтверждённый шаг; задним числом \enquote{додавить} ту же MIT
не~пытаются. Европейское регулирование и~правила карточных сетей здесь сходятся:
при удалённом оформлении мандата на~серию MIT SCA применяется
на~этапе создания мандата; произвольная серия повторов
в~требование SCA не~попадает~\cite{eba-mit-qa,mc-tpr}.

При отказе по~регулярному списанию действуют процедуры самих карточных сетей.
Visa для подписочных сценариев устанавливает отдельные требования
к~уведомлению держателя карты~--- например, при окончании пробного
или промопериода ТСП обязан предупредить клиента не~менее
чем за~7~календарных дней до~первого полного списания; для повторной
отправки отклонённой операции MIT по~разрешающим кодам отказа
(\textit{decline}) Visa ограничивает количество попыток
(с~25.05.2025~--- до~20 за~30~дней; штраф начисляется
с~21-й попытки)~\cite{visa-core-rules}. Mastercard для отказов
по~регулярным платежам ограничивает повторную отправку
(до~10~попыток за~24~часа, до~35~попыток за~30~дней) и~требует
учитывать MAC, который определяет допустимый интервал
и~возможность повтора~\cite{mc-tpr}. Лимиты заданы там, где
совокупный уровень успеха повторов перестаёт расти, а~нагрузка
на~авторизационную инфраструктуру~--- продолжает.

Коды MAC~24--30 задают минимальное время ожидания перед повтором:
MAC~24~--- 1~час, MAC~25~--- 24~часа, MAC~26~--- 2~дня,
MAC~27~--- 4~дня, MAC~28~--- 6~дней, MAC~29~--- 8~дней,
MAC~30~--- 10~дней. Код~24 рассчитан на~временное ограничение
баланса, которое снимется к~следующему расчётному дню, а~старшие
коды~--- на~ситуацию, когда держатель ждёт пополнения в~следующем
расчётном цикле. MAC~03 (\textit{Do not try again}) полностью
запрещает повторную отправку по~этим реквизитам, а~MAC~01
(\textit{New account information available}) указывает, что
реквизиты устарели и~ТСП должен получить обновлённые данные
через сервис Mastercard Automatic Billing Updater (ABU;
см.~раздел~\ref{sec:subscriptions-account-updater}) или
от~клиента~\cite{mc-tpr,mc-abu,tabapay-mac}. Mastercard отслеживает
повторные попытки после запрета и~включает ТСП в~мониторинговую
программу за~игнорирование MAC.

Visa использует схожую логику через Response Code и~Transaction
Advisor Code. Отказ с~кодом, требующим аутентификации, возвращает
клиента в~3DS-цикл; повторная MIT здесь не~запускается. Отказ
с~кодом, допускающим повтор, ограничен 20~попытками за~30~дней
для одного ТСП и~набора реквизитов~\cite{visa-core-rules}.
Дальнейшие попытки Visa расценивает как нарушение правил.

Классификатор MAC по~коду из~DE~48 (Additional Data~--- Private
Use) сводит ответ к~трём решениям: остановить повторы (MAC~03),
запросить обновление реквизитов через ABU или Visa Account Updater
(VAU; MAC~01) или повторить через предписанный интервал
(MAC~24--30). Незнакомый код трактуется как STOP~--- при таком
умолчании мониторинг карточной сети не~срабатывает из-за неучтённого
сценария. Без такого классификатора биллинг либо повторяет
слишком часто и~попадает под мониторинг карточной сети, либо прекращает
попытки слишком рано и~теряет подписчиков:

\begin{figure}[htbp]
  \centering
  \includefiguresvg[width=.8\linewidth]{assets/figures/ch29-subscriptions/decline-decision-tree}
  \caption{Дерево решений по~MAC: STOP (03), UPDATE (01), WAIT (24--30).}\label{fig:decline-decision-tree}
\end{figure}

\codefile{Python}{samples/ch29-mac-retry/mac.py}

\FloatBarrier
\needspace{4\baselineskip}
\section{Обновление реквизитов}\label{sec:subscriptions-account-updater}

Сервисы обновления реквизитов закрывают сценарий, когда клиент
не~менял подписку, но~эмитент перевыпустил карту. Старый PAN
больше не~подходит; ТСП продолжает списывать по~устаревшим
данным.

VAU и~ABU~--- сервисы соответственно Visa и~Mastercard, которые
передают ТСП обновлённые реквизиты от~участвующих эмитентов,
когда карта истекла или была
перевыпущена~\cite{mc-abu-privacy,visa-vau}. В~правилах Mastercard
для регулярных платежей участие в~ABU формализовано, с~исключениями
по~отдельным типам карт~\cite{mc-tpr}. Visa называет аналогичный
механизм credential lifecycle management. Обновление PAN и~сетевых
токенов возможно в~реальном времени прямо во~время
авторизации~\cite{visa-lifecycle-management}.

Пакетное обновление работает по~расписанию: ТСП или PSP
отправляет файл с~реквизитами, карточная сеть возвращает обновлённые данные
через несколько часов. Для подписок с~месячным циклом этого
достаточно, если файл отправляется заблаговременно. Обновление
в~реальном времени встроено в~сам авторизационный запрос: когда
реквизиты устарели, ответ содержит обновлённые данные, и~ТСП
сохраняет их сразу. Visa описывает это как Real-Time Account
Updater~\cite{visa-lifecycle-management}; Mastercard передаёт
обновления через Payment Account Management API, который снабжает
ABU данными от~эмитента по~мере перевыпуска~\cite{mc-tpr}.
С~01.04.2025 Mastercard ввёл дополнительный штраф Credential
Continuity Program~--- \$0,09 за~каждую транзакцию по~устаревшим
реквизитам, если ТСП не~подключён к~ABU или игнорирует
обновления; ставка выросла втрое от~прежних
\$0,03~\cite{mc-tpr}.

Покрытие зависит от~участия эмитента. Не~каждый банк подключён
к~ABU или VAU, не~каждый тип карты попадает в~процесс обновления.
Visa рекомендует подключать все идентификационные диапазоны
эмитента (Bank Identification Number, BIN), включая
перезагружаемые предоплаченные карты (reloadable prepaid),
но~неперезагружаемые предоплаченные карты и~продукты малых
эмитентов часто остаются вне сервиса. Обновление реквизитов
снижает технический отток, но~не~устраняет его. Для карт вне
покрытия единственный путь~--- попросить клиента ввести новые
данные.

С~сетевыми токенами механика устроена иначе: когда ТСП
хранит не~PAN, а~сетевой токен, при перевыпуске карты привязка
обновляется автоматически~--- токен остаётся прежним, исходный
PAN (underlying PAN, реальный номер карты, на~который сеть
отображает токен) обновляется на~стороне
карточной сети~\cite{visa-lifecycle-management}. Задача обновления
переносится от~ТСП к~карточной сети, а~пакетные файлы для
токенизированных реквизитов становятся не~нужны.

\begin{figure}[htbp]
  \centering
  \includefiguresvg[width=\linewidth]{assets/figures/ch29-subscriptions/credential-update-flow}
  \caption{Account Updater (VAU/ABU) и~сетевая токенизация: сравнение по~шести параметрам.}\label{fig:credential-update-flow}
\end{figure}

MIT уходит по~карте с~истёкшим сроком, а~сервис обновления
реквизитов (ABU/VAU) недоступен~--- пакетный файл не~обработался
или запрос в~реальном времени (\textit{real-time}) вернул ошибку.
Эмитент отклоняет по~коду~\texttt{54} (Expired Card). Единственный
путь~--- уведомить клиента и~попросить обновить реквизиты вручную.
Когда ТСП вместо этого повторяет тот же запрос, он тратит
попытки впустую и~рискует попасть в~мониторинг карточной сети
по~избыточным ретраям.

Другой случай: сетевой токен обновился при перевыпуске карты,
underlying PAN изменился на~стороне TSP, но~первая же MIT
по~обновлённому фондирующему PAN (Funding PAN, FPAN~--- реальный
номер карты, в~который сеть детокенизирует токен при авторизации)
получает отказ с~кодом~\texttt{62} (Restricted Card), потому что
новая карта выпущена с~ограничениями по~коду категории ТСП
(Merchant Category Code, MCC) или по~каналу CNP. Токен жив,
но~авторизация не~проходит; retry бесполезен, потому что
ограничение лежит на~стороне эмитента. Дальше нужно уведомить
клиента, запросить новый платёжный инструмент или, если доступно,
обратиться к~эмитенту через PSP для снятия ограничения.

Involuntary churn (отток из-за платёжных отказов, не~по~решению
клиента) по~индустриальным оценкам Recurly, Spreedly и~других
вендоров биллинга~--- 3--5\,\% подписочной базы в~месяц; успешность
ABU/VAU при обновлении реквизитов по~тем же оценкам~--- порядка
85--90\,\% для карт участвующих эмитентов. Это агрегированные
показатели по~портфелям подписочных продуктов; публичных данных
карточных сетей здесь нет. Сетевые токены повышают этот показатель, потому
что обновление происходит автоматически на~стороне TSP
без зависимости от~пакетного цикла.

\section{Биллинг-движки как продуктовая надстройка}\label{sec:subscription-billing-engines}

SCF, NTID, MIT, CIT, recurring против unscheduled~--- это
\textbf{протокол карточной сети}. Над ним работает \textbf{продуктовая
надстройка биллинг-движка}: выставление счетов (invoicing),
пропорциональное начисление при изменении тарифа (proration),
повторные попытки списания при отказе (dunning), отмена и~возврат,
налоговая обработка.

Число краевых случаев (\textit{edge cases}~--- изменение цикла,
перерасчёт скидок, состояния заморозки (\textit{freeze}), валютные
конверсии при подписке в~другой валюте, налоги по~юрисдикции
клиента) превышает то, что попадает в~первичную постановку
самописного биллинга. Под эти случаи сформировался отдельный
от~PSP рынок промышленных биллинг-движков.

Stripe Billing~\cite{stripe-billing}~--- встроенное в~Stripe
решение, плотно интегрированное с~их Payments и~налоговым сервисом
Stripe Tax; поддерживает модели по~потреблению (\textit{usage-based})
и~по~фиксированной цене (\textit{fixed-price}), повторные попытки
списания (\textit{dunning}), декларирует восстановление 55\,\%
неуспешных списаний в~среднем по~портфелю. Chargebee~\cite{chargebee-overview}
и~Recurly~\cite{recurly-overview}~--- независимые SaaS-платформы
с~интеграцией, не~привязанной к~одному PSP, нацелены на~сегмент
B2B~SaaS: Chargebee ставит на~гибкость моделей ценообразования
(\textit{pricing}), Recurly~--- на~восстановление выручки
(\textit{revenue recovery}) и~ML-настройку dunning.
Lago~\cite{lago-overview}~--- открытый исходный код (\textit{open
source}) под AGPL-3.0, версия для самостоятельного развёртывания
(\textit{self-hosted edition}) без~ограничений по~объёму;
позиционируется прямой альтернативой Stripe Billing.
Kill~Bill~\cite{killbill-overview}~--- открытая платформа
повторяющейся выручки (\textit{recurring revenue}) под Apache~2.0,
развивается с~2010~года; рассчитана на~сложные сценарии
с~собственной бизнес-логикой поверх ядра.

Движок выбирают по~трём критериям: тип модели ценообразования
(\textit{usage-based}, \textit{subscription} или гибрид),
требование к~самостоятельному развёртыванию (\textit{self-hosting};
открытый исходный код против SaaS) и~независимость от~конкретного
PSP. Последнее становится критичным, когда ТСП работает через
несколько эквайеров параллельно: встроенные в~один PSP движки эту
степень свободы теряют.

\section{Снижение оттока в~подписках}\label{sec:subscriptions-anti-churn}

Отток подписчиков по~платёжным отказам сокращается там, где
сохранён технический контекст согласия. Корректная цепочка NTID
показывает эмитенту, что запрос~--- продолжение известного
мандата: эмитент применяет менее строгую политику антифрода,
и~доля одобрений на~MIT с~корректным NTID выше, чем
на~\enquote{сиротских} запросах без~истории.

Первая CIT должна пройти и~зафиксировать согласие: идентификатор
исходной авторизации и~токен реквизита сохраняются на~уровне
подписки, а~не в~записи последней транзакции. Каждый последующий
биллинговый цикл~--- MIT с~корректными флагами SCF и~ссылкой
на~ту же NTID. При перевыпуске карты сетевой токен или ABU/VAU
закрывают обновление реквизитов автоматически; для~карт вне
покрытия и~отказов с~MAC~01 единственный путь~--- запросить
данные у~клиента вручную.

Когда эмитент возвращает отказ, требующий аутентификации, MIT
не~повторяется~--- клиент проходит новый цикл CIT и~даёт свежее
согласие. Для отказов с~MAC~24--30 повтор идёт по~предписанному
интервалу, для~MAC~03~--- останавливается. Игнорирование этих
сигналов даёт два одинаково плохих исхода: мониторинг карточной сети
за~избыточные ретраи или преждевременную отписку платёжеспособного
клиента.

По~картам с~истёкшим сроком или ограничениями эмитента отток
уходит туда, где техника бессильна: клиент даже не~знает, что
у~него отказ. Минимальный набор уведомлений по~подписке закрывает
этот пробел~--- дата следующего списания, информация
о~неуспешной попытке и~канал обновления реквизитов.

\vfill
\practicenote{%
  Цепочка NTID~--- технический эквивалент письменного согласия:
  без~неё каждое автосписание начинается с~нуля.
}
\vfill
